
\section{Method}

\subsection{Overview}

This section describes our reproduction of Nystr mformer from Xiong et al. 
(2021). We do not introduce novel components.

\textbf{Figure 4.1} shows the pipeline.

\begin{lstlisting}
   +----------------------------------------------+
   |  Input hidden\_states $\in$ R^{B $\times$ N $\times$ D}         |
   +--------------------+-------------------------+
                        |
                         
   +----------------------------------------------+
   |  Q = query(x), K = key(x), V = value(x)      |
   |  (with LoRA adapters)                         |
   +--------------------+-------------------------+
                        |
                         
   +----------------------------------------------+
   |  RoPE (if applicable)                         |
   |  GQA repeat\_kv (if applicable)                |
   +--------------------+-------------------------+
                        |
                         
   +----------------------------------------------+
   |  Branching:                                    |
   |  - if decoder: dense causal attention         |
   |  - if N $\le$ 128: dense attention                |
   |  - else: Nystr m branch                        |
   +--------------------+-------------------------+
                        |
                         
   +----------------------------------------------+
   |  out\_proj (if not nested)                     |
   +----------------------------------------------+
\end{lstlisting}

\subsection{Three Execution Paths}

Nystr mformer in our implementation has \textbf{three paths}:

\subsubsection{Causal Decoder Path}

For autoregressive decoders (GPT-2, Qwen3.5), Nystr m approximation is 
not directly applicable because landmark pooling averages across tokens 
including future ones. We fall back to \textbf{dense causal attention}.

\subsubsection{Short Sequence Fallback}

For $N \leq \max(128, 2m)$, we use dense attention. Reasons:

\begin{itemize}
\item Nystr m approximation is sub-optimal when $N$ is small.
\item Dense attention is memory-efficient for short sequences.

\end{itemize}
\subsubsection{Nystr m Path}

For $N > \max(128, 2m)$ and non-causal attention:

\[
\text{Output} = \text{softmax}(Q K\_m^\top) \cdot \text{pinv}(A\_{mm}) \cdot \text{softmax}(Q\_m K^\top) V
\tag{4.1}
\]

where $A_{mm} = \text{softmax}(Q_m K_m^\top / \sqrt{d_h})$.

\subsection{Landmark Selection}

We use \texttt{adaptive\_avg\_pool1d}:

\begin{lstlisting}
Q\_flat = Q.reshape(B * H, N, d).transpose(1, 2)   # (BH, d, N)
Q\_landmarks = F.adaptive\_avg\_pool1d(Q\_flat, m).transpose(1, 2).reshape(B, H, m, d)
\end{lstlisting}

Same for $K$.

\textbf{Properties:}
\begin{itemize}
\item Deterministic.
\item $\mathcal{O}(N d_h)$ cost.
\item Adjacent tokens share landmarks.

\end{itemize}
\subsection{The Three Softmax Matrices}

\[
\begin{aligned}
F\_1 &= \text{softmax}(Q K\_m^\top / \sqrt{d\_h}) \in \mathbb{R}^{N \times m} \\
F\_2 &= \text{softmax}(Q\_m K\_m^\top / \sqrt{d\_h}) \in \mathbb{R}^{m \times m} \\
F\_3 &= \text{softmax}(Q\_m K^\top / \sqrt{d\_h}) \in \mathbb{R}^{m \times N}
\end{aligned}
\tag{4.2}
\]

Then:

\[
\text{Output} = F\_1 \cdot \text{pinv}(F\_2) \cdot (F\_3 V)
\tag{4.3}
\]

\textbf{Associative order:} Compute $F_3 V$ first ($m \times d_h$), then pinv 
($m \times m$), then $F_1$ ($N \times m$). This ensures $\mathcal{O}(N m d_h)$ 
total.

\subsection{Iterative Pseudoinverse}

Nystr mformer uses Newton-Schulz iteration:

\begin{lstlisting}
def \_iterative\_pinv(A, iters=3):
    try:
        return torch.linalg.pinv(A.float()).to(A.dtype)
    except Exception:
        # Newton-Schulz fallback
        norm\_1 = A.abs().sum(dim=-2).amax(dim=-1, keepdim=True).unsqueeze(-1)
        norm\_inf = A.abs().sum(dim=-1).amax(dim=-1, keepdim=True).unsqueeze(-1)
        alpha = 1.0 / (norm\_1 * norm\_inf + 1e-5)
        V = alpha * A.transpose(-2, -1)
        I = torch.eye(A.size(-1), device=A.device, dtype=A.dtype).unsqueeze(0).unsqueeze(0)
        for \_ in range(iters):
            V = torch.matmul(V, 2.0 * I - torch.matmul(A, V))
        return V
\end{lstlisting}

\textbf{Note.} Our implementation first tries \texttt{torch.linalg.pinv} (SVD-based). 
If it fails (e.g., singularity), we fall back to Newton-Schulz.

\textbf{Difference from original.} The original Nystr mformer uses Newton-Schulz 
exclusively. We use direct pinv when possible, which is more accurate.

\subsection{Weight Transfer from Pretrained Attention}

\subsubsection{Architecture-Aware Extraction}

\begin{table}[H]
\centering
\small
\begin{tabular}{ll}
\toprule
\textbf{Architecture} & \textbf{Extraction} \\
\midrule
DistilBERT & \texttt{q\_lin, k\_lin, v\_lin, out\_lin} \\
RoBERTa & \texttt{self.query, self.key, self.value, output.dense} \\
GPT-2 & Split \texttt{c\_attn}; \texttt{c\_proj} for output \\
Qwen3.5 & \texttt{q\_proj, k\_proj, v\_proj, o\_proj} \\
\bottomrule
\end{tabular}
\end{table}

\subsubsection{Conv1D Transposition}

\begin{lstlisting}
if is\_conv1d:
    c\_w = c\_w.t()
\end{lstlisting}

\subsubsection{Nested Encoder Bypass}

\begin{lstlisting}
if self.is\_nested\_encoder:
    out = context   # Skip out\_proj
else:
    out = self.out\_proj(context)
\end{lstlisting}

\subsubsection{GQA}

\begin{lstlisting}
if self.kv\_repeat \textgreater{} 1:
    K = repeat\_kv(K, self.kv\_repeat)
    V = repeat\_kv(V, self.kv\_repeat)
\end{lstlisting}

\subsection{The Full Algorithm}

\textbf{Algorithm 1:} Nystr m attention forward
\begin{lstlisting}
Input:  Q, K, V $\in$ R^{B $\times$ H $\times$ N $\times$ d}
        valid\_mask $\in$ {0,1}^{B $\times$ N}
        num\_landmarks m
Output: context $\in$ R^{B $\times$ H $\times$ N $\times$ d}

1:  if is\_causal or is\_decoder:
2:      # Dense causal attention
3:      A $\leftarrow$ softmax(Q K^T / $\sqrt{}$d + causal\_mask + valid\_mask)
4:      context $\leftarrow$ A V
5:      return context
6:
7:  if N $\le$ max(128, 2m):
8:      # Dense attention
9:      A $\leftarrow$ softmax(Q K^T / $\sqrt{}$d + valid\_mask)
10:     context $\leftarrow$ A V
11:     return context
12:
13: # Nystr m branch
14: Q\_m $\leftarrow$ adaptive\_avg\_pool(Q, m)
15: K\_m $\leftarrow$ adaptive\_avg\_pool(K, m)
16:
17: F1 $\leftarrow$ softmax(Q K\_m^T / $\sqrt{}$d)         # (B, H, N, m)
18: F2 $\leftarrow$ softmax(Q\_m K\_m^T / $\sqrt{}$d)       # (B, H, m, m)
19: F3 $\leftarrow$ softmax(Q\_m K^T / $\sqrt{}$d)         # (B, H, m, N)
20:
21: F2\_pinv $\leftarrow$ iterative\_pinv(F2)       # (B, H, m, m)
22: step1 $\leftarrow$ F3 V                        # (B, H, m, d)
23: step2 $\leftarrow$ F2\_pinv step1               # (B, H, m, d)
24: context $\leftarrow$ F1 step2                  # (B, H, N, d)
25:
26: return context
\end{lstlisting}

\subsection{Complexity Analysis}

\begin{table}[H]
\centering
\small
\begin{tabular}{lll}
\toprule
\textbf{Path} & \textbf{Time} & \textbf{Memory} \\
\midrule
Dense causal & $\mathcal{O}(N^2 d_h)$ & $\mathcal{O}(N^2)$ \\
Dense (short) & $\mathcal{O}(N^2 d_h)$ & $\mathcal{O}(N^2)$ \\
Nystr m & $\mathcal{O}(N m d_h)$ & $\mathcal{O}(N m)$ \\
\bottomrule
\end{tabular}
\end{table}

\textbf{Nystr m speedup:}

\[
\frac{\text{Dense}}{\text{Nystr m}} = \frac{N^2 d\_h}{N m d\_h} = \frac{N}{m}
\tag{4.4}
\]

For $N = 4096$, $m = 32$: \textbf{128$\times$ speedup}.

\subsection{Theoretical Notes}

\subsubsection{Approximation Error}

\textbf{Claim.} If the attention matrix $A$ has rank $\leq m$, the Nystr m 
approximation is exact.

\textbf{Proof.} Nystr m reconstruction is exact for rank-$m$ matrices. $\square$

In practice, $A$ is approximately low-rank. Error depends on the tail of 
the singular value spectrum.

\subsubsection{Newton-Schulz Convergence}

\textbf{Proposition.} The iteration $V_{k+1} = V_k (2I - A V_k)$ with 
$V_0 = \alpha A^\top$ converges to $A^+$ if $\|I - \alpha A A^\top\|_2 < 1$.

\textbf{Proof.} Classical. $\square$

\subsubsection{Adaptive Pooling as Landmark Selection}

\textbf{Observation.} Adaptive pooling selects landmarks as segment averages. 
This is equivalent to sampling at stride $N / m$ with local averaging.

\textbf{Comparison to k-means:} Adaptive pooling is deterministic and faster. 
K-means would give better landmarks at higher cost.

\subsection{Summary}

Nystr mformer:
\begin{enumerate}
\item Approximates attention via landmarks.
\item Uses adaptive pooling for landmarks.
\item Uses Newton-Schulz pinv (or direct pinv).
\item Falls back to dense for short sequences and decoders.
\item Does \textbf{not} reduce parameter count or KV cache.
\item Reduces attention \textbf{compute}.
\end{enumerate}
\begin{lstlisting}

---
\end{lstlisting}